\section{Local rooted constructions}

We first record three small rooted constructions and a local contraction
lemma. They will also be used in the unrooted argument.

The rooted-star lemma is related to~\cite[Lemma~4.10]{DNR2026}, which
gives three spokes at one of five roots when $|V(H)|+|X|\le8$.
Here independence of $X$ allows all $k-1$ spokes and the bound
$|V(H)|+|X|\le2k-1$; the degree assumptions are stated below.

\begin{lemma}[Rooted star]\label{lem:star}
Let $Z\subseteq V(H)$ have order $k$, and let $X\subsetneq Z$ be
independent. Suppose that
\[
 |V(H)|+|X|\le 2k-1,\qquad
 d_H(v)\ge k-1\quad(v\notin X),\qquad
 d_H(x)\ge k-|X|\quad(x\in X).
\]
Then, for some $z\in Z\setminus X$, the set
$(V(H)\setminus Z)\cup\{z\}$ is connected and adjacent to every vertex
of $Z\setminus\{z\}$.
\end{lemma}

\begin{proof}
Put $T=V(H)\setminus Z$. A component $D$ of $H[T]$ misses at most
$|D|$ vertices of $Z$: any vertex of $D$ has at least $k-|D|$
neighbours in $Z$. Consequently at most $|T|$ roots miss some component
of $H[T]$. Since $|T|+|X|\le k-1$, there is a vertex
$z\in Z\setminus X$ adjacent to every component. Thus $T\cup\{z\}$
is connected, including when $T=\varnothing$.

A vertex $r\in Z\setminus\{z\}$ which has no neighbour in $T$ must
be adjacent to $z$. For $r\notin X$ this follows from $d_H(r)\ge k-1$;
for $r\in X$, independence of $X$ and $d_H(r)\ge k-|X|$ force
adjacency to every vertex of $Z\setminus X$.
\end{proof}

\begin{lemma}[Rooted triangle]\label{lem:triangle}
Let $(R,S)$ be internally three-connected, where $|S|=3$.
Suppose that $R-S$ is nonempty and connected, each root has a neighbour
outside $S$, and every nonroot has degree at least four. Then $R$
contains an $S$-rooted triangle minor.
\end{lemma}

\begin{proof}
Three distinct vertices $a,b,c$ of a two-connected graph have a rooted
triangle minor. To see this, take a cycle through $a,b$ and, if necessary,
a two-fan from $c$ to that cycle. Either their union contains a cycle
through all three vertices, or it contains a theta whose three paths
contain $a,b,c$ separately. In the latter case, assign one end of the
theta to the $a$-path and the other to the $b$-path; the interior of the
third path is the $c$-bag. These three bags are connected and pairwise
adjacent.

Consider the block-cut tree of $R$, representing a root which is a
cutvertex by its cutvertex node, and any other root by its block node.
If the median of the three root nodes is a block, the roots have three
distinct attachments to that block. The block is two-connected, so the
preceding observation supplies a triangle rooted at the attachments.
Appending the three disjoint exterior root paths gives the result.

If the median is a cutvertex $c$, each component of $R-c$ contains at
most one root. The nonroots in any such component would have boundary
contained in $c$ and that possible root, contrary to internal
three-connectivity. Hence $c$ is the only nonroot, but then its degree
is at most three. This is also impossible.
\end{proof}

We use the following terminal construction from the proof of
\cite[Lemma~3.1]{DNR2026}. Assume that four independent roots are the
leaves of a tree $T$. Its two degree-three branching vertices $v_1,v_2$
may coincide at a degree-four vertex. Choose $T$ to minimise the length
of the path between these vertices. If, after deleting $v_1,v_2$, every
partition of the roots into two pairs has a path joining the pairs,
then the graph has a rooted dart. This part of the cited proof uses
only the stated tree and path conditions, not a density hypothesis.

\begin{lemma}[Rooted dart]\label{lem:dart}
An internally four-connected graph with four prescribed roots, at least
two nonroots, and nonroot minimum degree at least five contains a rooted
dart.
\end{lemma}

\begin{proof}
Write $(F,Z)$ for the rooted graph and induct on $|V(F)|$. Delete the
root--root edges; this preserves all nonroot degrees and boundaries.
Every partition of $Z$ into two pairs has two disjoint paths between
the pairs. Otherwise Menger's theorem gives a separation of order at
most one, with one pair on each side. The nonroots on either open side
then have boundary contained in that pair and the separator, of order
at most three. Internal four-connectivity would put all nonroots in
the separator, contradicting their number.

Each component of $F-Z$ sees all four roots. It therefore supplies a
tree whose leaves are exactly $Z$. Choose such a tree with a shortest
branching path, with ends $v_1,v_2$. Unless $\{v_1,v_2\}$ separates
two roots from the other two, the terminal construction above gives a
dart. The case $v_1=v_2$ is covered, since a one-vertex separation of
this kind was already excluded.

In the remaining case let $(B_1,B_2)$ be the separation, with
$B_1\cap B_2=\{v_1,v_2\}$ and two roots in each open side. Put
\[
 Z_i=(Z\cap B_i)\cup\{v_1,v_2\},\qquad W_i=B_i\setminus Z_i.
\]
If both $W_i$ are empty, the degree bound makes $v_1,v_2$ adjacent to
each other and to every root. Use $v_1$ as the dart's helper and absorb
$v_2$ into one root bag: the resulting root star contains the required
three-vertex path.

Otherwise choose a nonempty $W_i$. It has at least two vertices, since
a single vertex would have at most the four neighbours in $Z_i$.
Every subset of $W_i$ has the same neighbours in $F[B_i]$ as in $F$.
Thus $(F[B_i],Z_i)$ is internally four-connected and retains the
nonroot degree bound. It omits the two opposite roots, so induction
provides a $Z_i$-rooted dart.

Take the two disjoint paths between the two original root pairs
established above. They cross the separator at different vertices.
Their segments after their last visits to $\{v_1,v_2\}$ lie outside
$B_i$, apart from their initial vertices, and end at the two opposite
roots. Append these segments to the bags at $v_1,v_2$. They are disjoint
from the other bags and from the helper, and restore the four original
roots without losing any contact.
\end{proof}

Call a graph \emph{quasi five-connected} if it is four-connected and has
no four-separation with at least two vertices on each open side. The
following local version of the crossing argument in
\cite[Lemma~1]{Kou2025} does not assume contraction-criticality.

\begin{lemma}\label{lem:quasi}
Let $Q$ be five-connected, and let $A$ have boundary $S=N_Q(A)$ of
order five. Suppose that $|A|\ge3$ and
$|V(Q)\setminus(A\cup S)|\ge2$. If $x\in S$ has a unique neighbour
$a$ in $A$, then $Q/xa$ is quasi five-connected.
\end{lemma}

\begin{proof}
An edge contraction preserves four-connectivity. A forbidden
four-separation of $Q/xa$ would lift to a five-cut $T$ of $Q$ containing
$x,a$, with open sides $B,D$ of order at least two. Write
$C=V(Q)\setminus(A\cup S)$ and
\[
 p=|A\cap T|,\quad r=|C\cap T|,\quad s=|S\cap T|,\quad
 u=|S\cap B|,\quad w=|S\cap D|.
\]
Then $p+r+s=u+w+s=5$ and $s\ge1$.

If $A\cap B\ne\varnothing$, its boundary is contained in
$(A\cap T)\cup(S\cap T)\cup(S\cap B)$ and omits $x$, whose unique
neighbour in $A$ lies in $T$. Five-connectivity gives $p+s+u\ge6$,
so $p>w$ and $u>r$. The opposite corner $C\cap D$ has boundary of
order at most $r+s+w<5$, and is therefore empty. Symmetrically,
$A\cap D\ne\varnothing$ implies $p>u$, $w>r$, and
$C\cap B=\varnothing$.

If both $A$-corners are nonempty, then $C\subseteq T$, so $r\ge2$,
whereas $u+w\ge2r+2\ge6>5-s$. If neither is nonempty, then
$p=|A|\ge3$ and $r\le1$. Some $C$-corner is nonempty, say $C\cap B$;
its boundary gives $u\ge p$. Hence $w\le1$ and $C\cap D$ is empty,
contrary to $|D|\ge2$.

Finally suppose only $A\cap B$ is nonempty. Then $D=S\cap D$, so
$w\ge2$. Since $p>w$, we have $r\le1$. But $C\cap B$ has boundary
of order at most $r+s+u=5+r-w\le4$, so it too is empty. This leaves
$|C|=r\le1$, the final contradiction.
\end{proof}

\section{Two rooted helpers}

For five prescribed roots $X$, say that a model has the
\emph{helper property} if it consists of five prescribed-root bags and
two further nonempty connected bags, all pairwise disjoint, and at least
ten of the eleven possible helper--root and helper--helper contacts
are present. The two helpers contain no root; no root--root contact is
required.

\begin{theorem}\label{thm:helper}
Every $4$-light five-rooted graph $(G,X)$ with $\rho_4(G,X)\ge2$ has
the helper property.
\end{theorem}

The proof first restricts a minimum counterexample. It then uses a
weighted neighbourhood to find an edge with few common neighbours.
An auxiliary clique allows Mader's atom theorem to be applied without
requiring intersections of positive-density fragments to remain positive.

\subsection{Minimum-counterexample reductions}

Suppose the theorem is false, and choose a counterexample minimising
$(|V(G)|,|E(G)|)$ lexicographically. Throughout this section, contractions
identify at most one original root in each preimage. A rooted model in
a minor lifts by replacing each of its vertices with the fixed connected
preimage; in particular, all five prescribed roots remain distinct.

We use the following consequence of the isolator argument in
\cite[Observation~4.1 and proof of Corollary~4.2]{DNR2026}. Suppose a
root-free region has a helper-property model on a five-vertex boundary,
or two helpers full to a larger boundary. Either five disjoint exterior
paths attach the original roots to five boundary bags, or a separator of
order $k\le4$ isolates the region and has a linkage to $k$ of its
boundary bags. A helper model restricted to those
$k$ bags contains a rooted $K_k$ for $k\le3$, and a rooted dart for
$k=4$. Indeed, retain a helper adjacent to all four roots and absorb the
other into a root bag to obtain two root contacts. For three roots,
absorb the helpers into two root bags; smaller cases follow by restriction.
The same conclusion holds when both helpers are full to a larger
boundary. If the isolated interior has at least two vertices and
nonpositive density, it is a reducible fragment. The linkage is outside
the original region, except for its ends, so it does not use a helper
vertex. This is the only form of the isolator argument needed below.

\begin{lemma}\label{rh:basic}
The minimum counterexample has the following properties:
\begin{enumerate}
\item $X$ is independent, and there is no reducible root-free fragment
with boundary of order at most four.
\item Every proper root-free set $Y$ with $|N_G(Y)|=5$ satisfies
$\rho_G(Y)\le1$.
\item $G-X$ is connected, every nonroot has at most three root neighbours,
and every root has degree at least two.
\item $\rho_4(G,X)=2$.
\end{enumerate}
\end{lemma}

\begin{proof}
Root--root edges contribute neither to the density nor to a required
helper contact, and cannot occur inside a bag containing just one root.
Thus minimality makes $X$ independent. A maximal reducible fragment
would give, by \cite[Lemma~3.3]{DNR2026}, a smaller $4$-light rooted
minor of nondecreasing density. Its helper model would lift, proving
(1).

For (2), put $S=N_G(Y)$. The induced rooted graph $(G[Y\cup S],S)$
is $4$-light, since its root-free sets retain their incident edges and
neighbours. If its density is at least two, it is a smaller graph with
a helper model by minimality. Five exterior root paths would extend
that model to $X$. Otherwise the isolator argument gives a reducible
fragment: its interior contains the two helpers and has nonpositive
density by $4$-lightness. Both possibilities contradict (1).

The densities of the components of $G-X$ sum to $\rho_4(G,X)$. Each
positive-density component sees all five roots, by $4$-lightness. Two
such components give two full helpers, with only their mutual contact
possibly missing. Hence there is exactly one positive component and
its density is at least two. Any other component would make this a
proper five-boundary set, contrary to (2). Thus $G-X$ is connected and
sees every root.

Suppose a nonroot $z$ has $d\ge4$ root neighbours. For the components
$Q_i$ of $G-(X\cup\{z\})$,
\[
 \sum_i\rho_G(Q_i)=\rho_4(G,X)+4-d.
\]
Each $Q_i$ sees $z$, and each positive one sees at least four roots.
If $d=5$, there is a positive component; it and $\{z\}$ are adjacent
helpers with at least ten contacts. If $d=4$, let $x$ be the missing
root. A component seeing all roots already gives the required model.
Otherwise each positive component has exactly four root neighbours,
and has density at most one by (2). Their sum is at least two, so there
are at least two positive components. Choose a component $R$ seeing
$x$, and a positive component $Q\ne R$. The helpers $\{z\}\cup R$
and $Q$ are adjacent and have respectively five and at least four root
contacts. This again contradicts the choice of $G$.

If a root $x$ had just one neighbour $v$, the side $G-x$, rooted at
$(X\setminus\{x\})\cup\{v\}$, would have density
$\rho_4(G,X)+4-|N_G(v)\cap X|\ge3$. By $4$-lightness, all five
proposed boundary vertices must be actual neighbours of its interior.
It is therefore a proper five-boundary side, contradicting (2). This
proves (3).

Finally suppose $\rho_4(G,X)\ge3$ and delete an edge with a nonroot
end. The result has density at least two but, by minimality, is not
$4$-light. A positive small-boundary set $Y$ becomes permissible only
because the deleted edge joins $Y$ to a vertex outside its closed
side. Restoring that edge gives a five-boundary set of density at least
two. By (2) its closed side is all of $G$, and its boundary is $X$.
The restored root endpoint then has only one neighbour, contrary to
(3). Therefore the density is exactly two.
\end{proof}

\begin{lemma}\label{rh:internal}
Every nonroot has degree at least five, and $(G,X)$ is internally
five-connected.
\end{lemma}

\begin{proof}
Deleting a nonroot $v$ of degree at most four leaves density
$2+4-d_G(v)\ge2$. If $G-v$ has a positive root-free set $Y$ with at
most four neighbours, restoring $v$ must add a neighbour to $Y$;
otherwise $Y$ already violates $4$-lightness in $G$. Its restored
density is at least two and its boundary has order five. This is a
proper side, since its boundary contains the nonroot $v$ and could not
also contain all five roots. Lemma~\ref{rh:basic}(2) excludes it.
Thus deletion stays in the induction class, a contradiction.

If internal five-connectivity fails, choose a connected root-free set
$Y$ of minimum boundary order $k\le4$, and put $S=N_G(Y)$. The rooted
shore $(G[Y\cup S],S)$ is internally $k$-connected, retains all nonroot
degrees, and has nonpositive density. For $k=0,1,2$, it contains a rooted
$K_k$ (using a path through $Y$ when $k=2$). For $k=3$, apply
Lemma~\ref{lem:triangle}. For $k=4$, the degree bound implies $|Y|\ge2$,
and Lemma~\ref{lem:dart} applies. In every case $Y$ is reducible,
contrary to Lemma~\ref{rh:basic}(1).
\end{proof}

Complete $X$ to a clique, obtaining $Q_0=G+K_X$. It is five-connected:
after deleting at most four vertices the surviving roots lie in one
clique, and another component would contradict internal five-connectivity.
For an original edge $e$ having a nonroot end, let $t(e)$ be its number
of common neighbours in $Q_0$. The rooted density after contraction is
\begin{equation}\label{rh:contract-density}
 \rho_4(G/e,X_e)=2+3-t(e),
\end{equation}
where $X_e$ denotes the unchanged or naturally identified roots. The
completion counts exactly the edges which become root--root edges when
one endpoint is a root.

\begin{lemma}\label{rh:blocker}
If $e=uv$ has a nonroot end and $t(e)\le3$, there is a proper root-free
set $Y$ such that
\[
 |N_G(Y)|=5,\qquad \rho_G(Y)=1,\qquad \{u,v\}\subseteq N_G(Y),
\]
and no vertex of $Y$ is adjacent to both $u$ and $v$.
\end{lemma}

\begin{proof}
By \eqref{rh:contract-density} and minimality, $G/e$ is not $4$-light.
Let $Y$ be a positive root-free set witnessing this. The identified
vertex cannot lie outside its closed side, because then $Y$ is a
root-free set of boundary at most four in $G$. Nor can it lie inside
$Y$: its preimage would have the same small boundary (and would be
root-free; a contracted root cannot be interior). It therefore lies
on the boundary. Splitting it gives exactly five neighbours in $G$.
The side is proper, since its boundary contains a nonroot and hence
cannot contain all five original roots. If $c$ vertices of $Y$ see
both endpoints, then
$\rho_G(Y)=\rho_{G/e}(Y)+c$. Lemma~\ref{rh:basic}(2) forces both
densities to equal one and $c=0$.
\end{proof}

\begin{lemma}\label{rh:five-edge}
There are no adjacent degree-five nonroots with exactly three common
neighbours.
\end{lemma}

\begin{proof}
Suppose $p,q$ are such vertices. Write their common neighbours as the
three-set $A$, and their respective exclusive neighbours as $u,v$.
Lemma~\ref{rh:blocker} applied to $pq$ gives a density-one set $T$
with boundary $\{p,q\}\cup D$, where $|D|=3$, avoiding $A$.
Since both $p,q$ have neighbours in $T$, their neighbourhoods force
$u,v\in T$; in particular these are nonroots.

Delete $p,q$ from this shore and reroot at $\{u,v\}\cup D$:
\[
 R=G[T\cup D],\qquad Z=\{u,v\}\cup D.
\]
The graph $(R,Z)$ is $4$-light. Indeed, every subset of its nonroots
$T\setminus\{u,v\}$ has exactly the same neighbours and incident
edges as in $G$, since only $u,v$ in $T$ see $p,q$.
Put $a_0=e_G(\{u,v\},D)$ and $\delta=1$ if $uv\in E(G)$, and
$\delta=0$ otherwise. Then
\[
 \rho_4(R,Z)
 =(4|T|+1-2-\delta-a_0)-4(|T|-2)
 =7-\delta-a_0.
\]
This is exactly the number $m\in\{0,\ldots,7\}$ of missing edges
from $uv$ together with the six edges between $\{u,v\}$ and $D$.
The rooted universality theorem \cite[Theorem~2.6]{DNR2026} supplies
all these missing edges simultaneously: its table contains every
$m$-edge graph in this range (the exceptional graph at density three
has four edges). Retain the root edges already present. We obtain
five rooted bags in $R$, with the bags at $u,v$ adjacent to one another
and to all three bags at $D$.

Append $p$ to the $u$-bag and $q$ to the $v$-bag. These are adjacent
root-free helpers, full to the boundary $A\cup D$. For $A\cap D$
use the existing $D$-bags; for $A\setminus D$ use singletons. These
choices are disjoint because $T\cap A=\varnothing$.

The root-free set $U=T\cup\{p,q\}$ has boundary exactly $A\cup D$.
If five exterior paths link $X$ to this boundary, retain the five bags
at their distinct ends and append the paths. This gives the helper
property rooted at $X$. Otherwise the isolator argument gives a
separator of order at most four, a rooted clique or dart on it, and a
root-free interior containing $U$. That interior has at least two
vertices and nonpositive density, so is reducible. Both alternatives
are impossible.
\end{proof}

\subsection{Eliminating degree five}

\begin{lemma}\label{rh:quasi-light}
If $u,v$ are nonroots and $Q_0/uv$ is quasi five-connected, then
$(G/uv,X)$ is $4$-light.
\end{lemma}

\begin{proof}
Root completion does not change any root-free boundary or incident
density. Every root has degree at least five in $Q_0/uv$: its at least
two nonroot neighbours merge to at least one, and it retains the four
other roots as neighbours. A root-free set with boundary at most four
has nonempty opposite side, because there are five roots. Its boundary
therefore has order four, and quasi five-connectivity makes one open
side a singleton. A singleton root-free side has degree four and
density zero. A singleton opposite side would be a root of degree at
most four, which is impossible.
\end{proof}

\begin{lemma}\label{rh:degree-six}
Every nonroot of $G$ has degree at least six.
\end{lemma}

\begin{proof}
Suppose a nonroot $v$ has degree five. First assume $Q_0[N_G(v)]$ is
not complete, and choose a nonedge $ab$ in it. Its ends are not both
roots. The graph
\[
 H=G-v+ab
\]
has rooted density two and is a proper rooted minor: contract $va$
and delete the unwanted new edges. If $H$ is $4$-light, its helper
model lifts, contrary to minimality.

Otherwise take a positive root-free set $Y$ of $H$ with at most four
neighbours. Put $k=e_G(v,Y)$ and let $c$ indicate whether the added
edge $ab$ is incident with $Y$. Then
\[
 N_G(Y)\subseteq N_H(Y)\cup\{v\},\qquad
 \rho_G(Y)=\rho_H(Y)+k-c.
\]
Internal five-connectivity gives $|N_G(Y)|=5$ and $k\ge1$. This side
is proper, since its boundary contains the nonroot $v$. Consequently
\[
 1\ge\rho_G(Y)=\rho_H(Y)+k-c\ge1+k-c,
\]
forcing $\rho_G(Y)=\rho_H(Y)=k=c=1$. Exactly one endpoint of $ab$
lies in $Y$; call it $a$. Then $a$ is a nonroot and the unique neighbour
of $v$ in $Y$. The other endpoint $b$ is among the four vertices of
$N_G(Y)\setminus\{v\}$, since the added edge cannot introduce a fifth
neighbour in $H$.

The set $Y$ is not a singleton, which would make $ab$ an edge of $G$.
If $|Y|=2$, its density one and the nonroot degree bound force adjacent
degree-five vertices with three common neighbours, contrary to
Lemma~\ref{rh:five-edge}. Hence $|Y|\ge3$. Its opposite side in
$Q_0$ has at least two vertices: a singleton would be an original root
$z$, with boundary $(X\setminus\{z\})\cup\{v\}$, and independence
of $X$ would give $d_G(z)\le1$.

Lemma~\ref{lem:quasi} now applies to the unique edge $va$ into $Y$.
Thus $Q_0/va$ is quasi five-connected, and Lemma~\ref{rh:quasi-light}
makes $G/va$ $4$-light. Since $v$ has degree five and $b$ misses $a$,
the edge $va$ has at most three common neighbours. Its contraction
retains rooted density at least two and decreases order, a contradiction.

It remains that $Q_0[N_G(v)]$ is a clique. Internal five-connectivity
and Menger's theorem give five disjoint paths from $X$ to $N_G(v)$,
with interiors outside $N_G[v]$. Roots already in $N_G(v)$ use trivial
paths. Choose $z\in N_G(v)\setminus X$, which exists because $v$ has
at most three root neighbours. Use $\{v\},\{z\}$ as helpers, and the
five linkage paths as root bags, removing $z$ from its own path. That
trimmed bag is nonempty because $z$ is not an original root. The helper
$z$ meets all five root bags, $v$ meets the other four, and $vz$ joins
the helpers. Every used clique edge has nonroot end $z$, so is an actual
edge of $G$. This gives the ten required contacts.
\end{proof}

\subsection{Weighted neighbourhoods and an atom}

Let $E_0$ be the original edges with a nonroot end and at most three
common neighbours in $Q_0$. Each edge of $E_0$ has a five-boundary
set as in Lemma~\ref{rh:blocker}.

\begin{lemma}\label{rh:weighted}
If $v$ is a nonroot and $d_G(v)+|N_G(v)\cap X|\le9$, then some edge
of $E_0$ is incident with $v$.
\end{lemma}

\begin{proof}
Suppose instead that $t(vu)\ge4$ for every neighbour $u$ of $v$.
Put $H=G[N_G(v)]$, $X_0=N_G(v)\cap X$, and $r=|X_0|$. Vertices of
$H-X_0$ have degree at least four in $H$; vertices of $X_0$ have degree
at least $5-r$, since root completion adds exactly $r-1$ common
neighbours to such an edge. Also $X_0$ is independent and
$|V(H)|+r\le9$.

There are five disjoint paths from $X$ to $N_G(v)$ in $G-v$.
Otherwise a separator of order at most four would isolate, in $G$,
a root-free component containing $v$. Truncate the paths at their first
entries to $N_G(v)$. Their five ends form a set $Z$ containing $X_0$,
and their interiors avoid the neighbourhood. Lemma~\ref{lem:star}
gives a central bag
\[
 C=(N_G(v)\setminus Z)\cup\{z\},\qquad z\in Z\setminus X_0,
\]
adjacent to the four other ends. Use $C$ and $\{v\}$ as helpers,
keeping the four other linkage paths and trimming $z$ from its path.
The trimmed root bag is nonempty; all bags are disjoint. The helper
$C$ has five root contacts, $v$ has four, and they are adjacent. This
is the helper property, a contradiction.
\end{proof}

In particular, $E_0$ is nonempty. If $m=|V(G)\setminus X|$, the
nonroot degree bound and at most three root neighbours give $m\ge4$,
while
\begin{equation}\label{rh:total-weight}
 \sum_{v\notin X}\bigl(d_G(v)+|N_G(v)\cap X|\bigr)=8m+4.
\end{equation}
Some nonroot therefore has weight at most nine.

We use Mader's atom theorem in the form stated in
\cite[Theorem~5]{KS2017}. In a graph $Q$ of connectivity $k$, let
$\mathcal S$ be any family of vertex sets. An $\mathcal S$-fragment
is a nonempty proper union of components of $Q-T$, where $T$ is a
$k$-cut containing a member of $\mathcal S$. An $\mathcal S$-atom
is such a fragment of minimum order. For an atom $A$, put
$T_A=N_Q(A)$ and $\overline A=V(Q)\setminus(A\cup T_A)$. If some
$S\in\mathcal S$ and $k$-cut $T$ satisfy
\begin{equation}\label{rh:atom-hyp}
 S\subseteq T\setminus\overline A,\qquad A\cap T\ne\varnothing,
\end{equation}
then
\begin{equation}\label{rh:atom-conclusion}
 A\subseteq T,\qquad |A|\le\tfrac12|T\setminus T_A|.
\end{equation}
There is no assumption that every edge is noncontractible.

\begin{proof}[Completion of the proof of Theorem~\ref{thm:helper}]
Add a set $W$ of $|V(G)|+1$ vertices and make $X\cup W$ a clique,
with no edges from $W$ to $V(G)\setminus X$; call the resulting graph
$Q$. Internal five-connectivity shows that $Q$ is five-connected, and
$X$ separates $G-X$ from $W$, so its connectivity is exactly five.
This padding changes neither boundaries of original root-free sets nor
common-neighbour counts of original edges with a nonroot end.

Let $\mathcal S=\{V(e):e\in E_0\}$. By Lemma~\ref{rh:blocker},
$\mathcal S$-fragments exist in $Q$, with an example of order at most
$m=|V(G)\setminus X|$. Choose an $\mathcal S$-atom $A$, among all
such fragments, so $|A|\le m$. It avoids $X\cup W$: otherwise its
five-boundary would force it to contain the whole surviving clique,
of order at least $|W|>m$. Thus $A\subseteq V(G)\setminus X$, and
its boundary is the same in $G$ and $Q$.

Its closed side is proper in $G$. If not, its five-boundary would be
exactly $X$, and could not contain the ends of an edge in $E_0$.
Lemma~\ref{rh:basic}(2) therefore gives $\rho_G(A)\le1$. The
nonroot degree bound excludes $|A|=1,2$, since for these orders
\[
 \rho_G(A)\ge2|A|-\binom{|A|}{2}\ge2.
\]
Writing $S_A=N_G(A)$, double-counting incident edges gives
\[
 \sum_{v\in A}\bigl(d_G(v)+|N_G(v)\cap S_A|\bigr)
   =8|A|+2\rho_G(A)\le8|A|+2.
\]
As $|A|\ge3$, some vertex $v\in A$ has the displayed weight at most
eight. Its root weight is no larger, so Lemma~\ref{rh:weighted}
supplies an incident edge $f\in E_0$. Both ends of $f$ lie in
$A\cup T_A$.

Apply Lemma~\ref{rh:blocker} to $f$, and let $T$ be the boundary of
its resulting fragment. In $Q$, this is a five-cut containing both
ends of $f$ and meeting $A$ at $v$. Consequently
$V(f)\subseteq T\setminus\overline A$, so \eqref{rh:atom-hyp} holds.
Mader's theorem yields $|A|\le5/2$, contradicting $|A|\ge3$.

The auxiliary clique has only selected an atom: no model has been
constructed in the padded graph. The contradiction proves the theorem
in the original rooted graph.
\end{proof}
